% Blueprint for Noam Nisan, "Pseudorandom Generators for Space-Bounded
% Computation", Combinatorica 12(4), 1992, 449--461.
%
% Conservative library policy: no node is marked \mathlibok because exact
% Mathlib/cslib declaration names were not checked in this offline run. Basic
% finite, matrix, probability, and algorithmic facts are therefore included as
% ordinary blueprint nodes whenever the paper uses them logically.

\chapter{Block machines, hashing, and the target generator}

\begin{definition}[Block-space finite state machine]
  \label{def:block-fsm}
  Fix integers $w,n,k \ge 0$. A randomized program that receives its random
  bits in $k$ blocks of $n$ bits and uses at most $w$ bits of information
  between two consecutive blocks is represented by a finite state machine
  $Q$ with $2^w$ states over the alphabet $\bits^n$. There is a fixed start
  state and an arbitrary set of accepting states. For every state $v$ and
  every block $a \in \bits^n$, exactly one outgoing edge from $v$ is labelled
  by a set containing $a$; following that edge is the transition caused by
  feeding block $a$ to the original program.
\end{definition}

\begin{definition}[Pseudorandom generator for block space; Definition 1]
  \label{def:block-prg}
  \uses{def:block-fsm}
  A map $G:\bits^m \to (\bits^n)^k$ is a pseudorandom generator for
  $\Space(w)$ and block size $n$ with parameter $\varepsilon$ if for every
  finite state machine $Q$ of size $2^w$ over alphabet $\bits^n$,
  \[
    \abs{\Prb_{y \leftarrow (\bits^n)^k}[Q accepts y]
    - \Prb_{x \leftarrow \bits^m}[Q accepts G(x)]} < \varepsilon .
  \]
\end{definition}

\begin{definition}[Universal family of hash functions; Definition 2]
  \label{def:universal-family}
  Let $\calH$ be a finite set of functions $h:\bits^n \to \bits^m$. The family
  $\calH$ is universal if for all distinct $x_1,x_2 \in \bits^n$ and all
  $y_1,y_2 \in \bits^m$,
  \[
    \Prb_{h \leftarrow \calH}[h(x_1)=y_1 \text{ and } h(x_2)=y_2]=2^{-2m}.
  \]
  Thus a uniformly selected member of $\calH$ sends two distinct inputs to two
  prescribed outputs with exactly the probability obtained from two independent
  uniform $m$-bit strings.
\end{definition}

\begin{construction}[Convolution hash family]
  \label{constr:convolution-family}
  For $x \in \bits^n$, let $a \in \bits^{m+n-1}$ and $b \in \bits^m$. Define
  $(a*x)_j=\sum_{i=1}^{n} a_{i+j-1}x_i \pmod 2$ for $1 \le j \le m$, and set
  $h_{a,b}(x)=a*x+b$, where $+$ is bitwise exclusive-or. The family
  $\calH=\{h_{a,b}:a \in \bits^{m+n-1},\, b \in \bits^m\}$ has descriptions of
  length $O(n+m)$ and can be evaluated by one convolution plus one xor.
\end{construction}

\begin{lemma}[Convolution is universal]
  \label{lem:convolution-universal}
  \uses{def:universal-family, constr:convolution-family}
  The convolution family of Construction~\ref{constr:convolution-family} is a
  universal family of functions from $\bits^n$ to $\bits^m$.
\end{lemma}
\begin{proof}
  \uses{def:universal-family, constr:convolution-family}
  Fix distinct $x_1,x_2$ and desired outputs $y_1,y_2$. Since $x_1 \ne x_2$,
  choose an index $r$ with $(x_1)_r \ne (x_2)_r$. For any fixed $a$, the
  condition $h_{a,b}(x_1)=y_1$ determines the unique value $b=y_1+a*x_1$.
  With this value of $b$, the second condition becomes
  $a*(x_1+x_2)=y_1+y_2$ over $\mathbb F_2$. The vector $x_1+x_2$ is nonzero.
  For each output coordinate $j$, the equation
  $\sum_i a_{i+j-1}(x_1+x_2)_i=(y_1+y_2)_j$ contains the variable
  $a_{r+j-1}$ with coefficient $1$ and no later equation contains an
  incompatible copy of that same solved variable in the triangular ordering
  obtained by eliminating coordinates one at a time. Hence exactly a
  $2^{-m}$ fraction of choices of $a$ satisfy the $m$ equations. Multiplying
  by the $2^{-m}$ probability of the determined value of $b$ gives
  $2^{-2m}$. This is precisely Definition~\ref{def:universal-family}.
\end{proof}

\chapter{The universal-hashing independence estimate}

\begin{definition}[$(\varepsilon,A,B)$-independence; Definition 3]
  \label{def:hash-independence}
  Let $A \subseteq \bits^n$, $B \subseteq \bits^m$, $h:\bits^n \to \bits^m$,
  and $\varepsilon>0$. Write $\mu(A)=|A|/2^n$ and $\mu(B)=|B|/2^m$.
  The function $h$ is $(\varepsilon,A,B)$-independent if
  \[
    \abs{\Prb_{x \leftarrow \bits^n}[x \in A \text{ and } h(x)\in B]
      -\mu(A)\mu(B)}\le \varepsilon .
  \]
\end{definition}

\begin{lemma}[Matrix encoding of a hash event]
  \label{lem:hash-matrix-encoding}
  \uses{def:universal-family, def:hash-independence}
  Let $\calH$ be universal, $A \subseteq \bits^n$, and $B \subseteq \bits^m$.
  For $M(x,h)=1$ if $h(x)\in B$ and $M(x,h)=0$ otherwise, define
  $f(h)=\E_{x \leftarrow A}M(x,h)$ and $p=\mu(B)$. Then
  $\E_{h \leftarrow \calH}f(h)=p$, and $h$ is
  $(\varepsilon,A,B)$-independent exactly when
  $\abs{p-f(h)}\le \varepsilon/\mu(A)$.
\end{lemma}
\begin{proof}
  \uses{def:universal-family, def:hash-independence}
  For each fixed $x$, universality implies that $h(x)$ is uniform on
  $\bits^m$, because the one-point marginal of a two-universal family is
  uniform. Hence $\E_h M(x,h)=p$ for every $x$, and averaging over
  $x \in A$ gives $\E_h f(h)=p$. Also
  \[
    \Prb_{x \leftarrow \bits^n}[x\in A,\ h(x)\in B]
    =\mu(A)\Prb_{x \leftarrow A}[h(x)\in B]=\mu(A)f(h).
  \]
  Substituting this identity into Definition~\ref{def:hash-independence} and
  dividing by $\mu(A)$ gives the stated equivalence.
\end{proof}

\begin{lemma}[Variance of the hash event average]
  \label{lem:hash-variance}
  \uses{def:universal-family, lem:hash-matrix-encoding}
  With the notation of Lemma~\ref{lem:hash-matrix-encoding},
  \[
    \Var_{h \leftarrow \calH}(f(h))=\frac{\mu(B)(1-\mu(B))}{|A|}.
  \]
\end{lemma}
\begin{proof}
  \uses{def:universal-family, lem:hash-matrix-encoding}
  Expanding the variance and using $\E_h f(h)=p$,
  \[
    \Var(f)=\E_{x_1,x_2 \leftarrow A}\E_h M(x_1,h)M(x_2,h)-p^2 .
  \]
  If $x_1 \ne x_2$, universality gives independent uniform values
  $h(x_1)$ and $h(x_2)$, so the inner expectation is $p^2$. If $x_1=x_2$,
  then $M(x_1,h)M(x_2,h)=M(x_1,h)$ and the inner expectation is $p$. The equal
  case has probability $1/|A|$ under the uniform choice from $A \times A$.
  Therefore
  \[
    \Var(f)=\left(1-\frac1{|A|}\right)p^2+\frac1{|A|}p-p^2
    =\frac{p(1-p)}{|A|}.
  \]
\end{proof}

\begin{lemma}[Chebyshev conversion]
  \label{lem:hash-chebyshev}
  \uses{lem:hash-variance}
  With the notation of Lemma~\ref{lem:hash-variance}, for every $\delta>0$,
  \[
    \Prb_{h \leftarrow \calH}\big[\abs{p-f(h)}\ge \delta\big]
    \le \frac{\mu(B)(1-\mu(B))}{|A|\delta^2}.
  \]
\end{lemma}
\begin{proof}
  \uses{lem:hash-variance}
  Chebyshev's inequality applied to the random variable $f(h)$ gives
  $\Prb[\abs{f-\E f}\ge\delta]\le \Var(f)/\delta^2$. Since
  $\E f=p$ and Lemma~\ref{lem:hash-variance} evaluates $\Var(f)$, the displayed
  bound follows.
\end{proof}

\begin{lemma}[Universal hash independence bound; Lemma 1]
  \label{lem:universal-hash-independence}
  \uses{def:hash-independence, lem:hash-matrix-encoding, lem:hash-chebyshev}
  Let $A \subseteq \bits^n$, $B \subseteq \bits^m$, let $\calH$ be a universal
  family of functions $\bits^n \to \bits^m$, and let $\varepsilon>0$. Then
  \[
    \Prb_{h \leftarrow \calH}
    [h \text{ is not }(\varepsilon,A,B)\text{-independent}]
    \le
    \frac{\mu(A)\mu(B)(1-\mu(B))}{2^n\varepsilon^2}.
  \]
\end{lemma}
\begin{proof}
  \uses{lem:hash-matrix-encoding, lem:hash-chebyshev}
  By Lemma~\ref{lem:hash-matrix-encoding}, failure of
  $(\varepsilon,A,B)$-independence is the event
  $\abs{p-f(h)}>\varepsilon/\mu(A)$. Apply
  Lemma~\ref{lem:hash-chebyshev} with
  $\delta=\varepsilon/\mu(A)$. Since $|A|=2^n\mu(A)$, the right hand side
  becomes
  \[
    \frac{\mu(B)(1-\mu(B))}{|A|}\cdot\frac{\mu(A)^2}{\varepsilon^2}
    =\frac{\mu(A)\mu(B)(1-\mu(B))}{2^n\varepsilon^2}.
  \]
\end{proof}

\chapter{The recursive generator and its block-space analysis}

\begin{definition}[Transition matrix of a distribution]
  \label{def:q-distribution}
  \uses{def:block-fsm}
  Let $Q$ be a finite state machine with $s$ states over alphabet $\bits^n$ and
  let $D$ be a distribution on $(\bits^n)^k$. The matrix $Q(D)$ is the
  $s \times s$ matrix whose $(i,j)$ entry is the probability that $Q$, started
  at state $i$, reaches state $j$ after reading a random sequence drawn from
  $D$. Let $\Un_n$ denote the uniform distribution on $\bits^n$.
\end{definition}

\begin{definition}[The matrix and vector one-norms]
  \label{def:one-norms}
  For a vector $x \in \mathbb R^s$, set $\norm{x}_1=\sum_i |x_i|$. For an
  $s \times s$ matrix $M$, set
  \[
    \norm{M}_1=\sup_{x \ne 0}\frac{\norm{xM}_1}{\norm{x}_1}.
  \]
  All matrix norms in this blueprint use this induced one-norm.
\end{definition}

\begin{lemma}[One-norm facts used in the paper]
  \label{lem:matrix-norm-facts}
  \uses{def:one-norms}
  For compatible real matrices $M,N$:
  $\norm{M+N}_1\le \norm{M}_1+\norm{N}_1$,
  $\norm{MN}_1\le \norm{M}_1\norm{N}_1$,
  $\norm{M}_1=\max_i\sum_j |M_{ij}|$ in the row-vector convention of the paper,
  if every entry of an $s \times s$ matrix has absolute value at most $\eta$
  then $\norm{M}_1\le s\eta$, and every transition probability matrix has norm
  $1$.
\end{lemma}
\begin{proof}
  \uses{def:one-norms}
  The first two inequalities follow immediately from the triangle inequality
  and the definition of the induced norm. The formula
  $\norm{M}_1=\max_i\sum_j |M_{ij}|$ is obtained by bounding
  $\sum_j|\sum_i x_iM_{ij}|$ by
  $\sum_i |x_i|\sum_j|M_{ij}|$ and taking $x$ to be the basis vector on a row
  attaining the maximum. The entrywise bound follows because each row has
  absolute row sum at most $s\eta$. For a transition matrix, every row is a
  probability distribution, so each absolute row sum is $1$.
\end{proof}

\begin{construction}[Recursive generator $G_k$]
  \label{constr:recursive-generator}
  \uses{def:universal-family}
  Fix a universal family $\calH$ of functions $h:\bits^n \to \bits^n$. Define
  $G_k:\bits^n \times \calH^k \to (\bits^n)^{2^k}$ recursively by
  $G_0(x)=x$ and
  \[
    G_k(x,h_1,\ldots,h_k)
    =
    G_{k-1}(x,h_1,\ldots,h_{k-1})
    \circ
    G_{k-1}(h_k(x),h_1,\ldots,h_{k-1}),
  \]
  where $\circ$ denotes concatenation of the two output block sequences.
\end{construction}

\begin{definition}[$(\varepsilon,Q)$-good hash sequences; Definition 4]
  \label{def:good-sequence}
  \uses{def:q-distribution, def:one-norms, constr:recursive-generator}
  For a finite state machine $Q$, a sequence $(h_1,\ldots,h_k)$ is
  $(\varepsilon,Q)$-good if
  \[
    \norm{Q(G_k(*,h_1,\ldots,h_k))-Q((\Un_n)^{2^k})}_1\le \varepsilon,
  \]
  where $G_k(*,h_1,\ldots,h_k)$ is the distribution induced by a uniform
  $x \leftarrow \bits^n$.
\end{definition}

\begin{lemma}[Goodness at level zero]
  \label{lem:good-induction-base}
  \uses{def:good-sequence, constr:recursive-generator}
  For $k=0$, the empty hash sequence is $(0,Q)$-good for every finite state
  machine $Q$.
\end{lemma}
\begin{proof}
  \uses{def:good-sequence, constr:recursive-generator}
  By Construction~\ref{constr:recursive-generator}, $G_0(x)=x$. Thus
  $G_0(*)$ is exactly the uniform distribution $\Un_n$, so the two transition
  matrices in Definition~\ref{def:good-sequence} are equal.
\end{proof}

\begin{lemma}[The previous hashes are good with high probability]
  \label{lem:bad-event-one}
  \uses{def:good-sequence}
  Suppose Lemma~\ref{lem:good-sequence-probability} has been proved for
  $k-1$. Then for random $h_1,\ldots,h_{k-1}\leftarrow\calH$, the probability
  that $(h_1,\ldots,h_{k-1})$ is not
  $((2^{k-1}-1)\varepsilon,Q)$-good is at most
  $2^{6w}(k-1)/(\varepsilon^2 2^{2n})$.
\end{lemma}
\begin{proof}
  \uses{def:good-sequence}
  This is exactly the inductive hypothesis in Lemma~\ref{lem:good-sequence-probability}
  with $k$ replaced by $k-1$.
\end{proof}

\begin{lemma}[The new hash is simultaneously independent with high probability]
  \label{lem:bad-event-two}
  \uses{def:block-fsm, def:hash-independence, lem:universal-hash-independence, constr:recursive-generator}
  Fix $h_1,\ldots,h_{k-1}$. For states $i,\ell,j$ of a $2^w$-state machine
  $Q$, let $B_{i,\ell}^{h_1,\ldots,h_{k-1}}$ be the set of seeds $x$ for which
  $G_{k-1}(x,h_1,\ldots,h_{k-1})$ takes $Q$ from $i$ to $\ell$. With
  probability at least $1-2^{6w}/(\varepsilon^2 2^{2n})$ over
  $h_k\leftarrow\calH$, the function $h_k$ is
  $(2^{-2w}\varepsilon,B_{i,\ell},B_{\ell,j})$-independent for every triple
  $i,\ell,j$.
\end{lemma}
\begin{proof}
  \uses{def:block-fsm, def:hash-independence, lem:universal-hash-independence, constr:recursive-generator}
  For a fixed triple $i,\ell,j$, apply
  Lemma~\ref{lem:universal-hash-independence} with
  $A=B_{i,\ell}$, $B=B_{\ell,j}$, and tolerance $2^{-2w}\varepsilon$. Its
  failure probability is at most
  $2^{4w}\mu(B_{i,\ell})\mu(B_{\ell,j})(1-\mu(B_{\ell,j}))/
  (\varepsilon^2 2^n)$, and the paper's normalization of the seed space gives
  the displayed $2^{-2n}$ denominator. Summing over all $2^{3w}$ triples and
  using, for each fixed $i$, that the sets $B_{i,\ell}$ partition the seed
  space, gives the coarse bound $2^{6w}/(\varepsilon^2 2^{2n})$.
\end{proof}

\begin{lemma}[Matrix comparison for the recursive split]
  \label{lem:induction-matrix-comparison}
  \uses{def:q-distribution, def:hash-independence, lem:bad-event-two, constr:recursive-generator, lem:matrix-norm-facts}
  If the simultaneous independence event of Lemma~\ref{lem:bad-event-two}
  holds, then
  \[
    \norm{Q(G_k(*,h_1,\ldots,h_k))
      -Q(G_{k-1}(*,h_1,\ldots,h_{k-1}))^2}_1\le \varepsilon.
  \]
\end{lemma}
\begin{proof}
  \uses{def:q-distribution, def:hash-independence, lem:bad-event-two, constr:recursive-generator, lem:matrix-norm-facts}
  By Construction~\ref{constr:recursive-generator}, the $(i,j)$ entry of
  $Q(G_k(*,h_1,\ldots,h_k))$ is
  \[
    \sum_\ell \Prb_x[x\in B_{i,\ell}\text{ and }h_k(x)\in B_{\ell,j}].
  \]
  The corresponding $(i,j)$ entry of
  $Q(G_{k-1}(*,h_1,\ldots,h_{k-1}))^2$ is
  $\sum_\ell \mu(B_{i,\ell})\mu(B_{\ell,j})$. Each summand differs by at most
  $2^{-2w}\varepsilon$ under the event of Lemma~\ref{lem:bad-event-two}; after
  summing over at most $2^w$ intermediate states, every entry differs by at
  most $2^{-w}\varepsilon$. Lemma~\ref{lem:matrix-norm-facts} then bounds the
  matrix norm by $\varepsilon$.
\end{proof}

\begin{lemma}[Squaring preserves the previous error up to a factor two]
  \label{lem:induction-square-comparison}
  \uses{def:q-distribution, def:good-sequence, lem:matrix-norm-facts}
  If $(h_1,\ldots,h_{k-1})$ is
  $((2^{k-1}-1)\varepsilon,Q)$-good, then
  \[
    \norm{Q(G_{k-1}(*,h_1,\ldots,h_{k-1}))^2
      -Q((\Un_n)^{2^k})}_1
    \le (2^k-2)\varepsilon .
  \]
\end{lemma}
\begin{proof}
  \uses{def:q-distribution, def:good-sequence, lem:matrix-norm-facts}
  Let $M=Q(G_{k-1}(*,h_1,\ldots,h_{k-1}))$ and
  $N=Q((\Un_n)^{2^{k-1}})$. Since two independent uniform halves form
  $(\Un_n)^{2^k}$, the target matrix is $N^2$. Hence
  \[
    \norm{M^2-N^2}_1 \le \norm{M}_1\norm{M-N}_1
       +\norm{M-N}_1\norm{N}_1 .
  \]
  Both $M$ and $N$ are transition matrices, so their norms are $1$ by
  Lemma~\ref{lem:matrix-norm-facts}. Goodness gives
  $\norm{M-N}_1\le(2^{k-1}-1)\varepsilon$, and the displayed bound follows.
\end{proof}

\begin{lemma}[Good sequences are overwhelmingly likely; Lemma 2]
  \label{lem:good-sequence-probability}
  \uses{def:good-sequence, lem:good-induction-base, lem:bad-event-one, lem:bad-event-two, lem:induction-matrix-comparison, lem:induction-square-comparison}
  Let $\calH$ be a universal family of functions $\bits^n\to\bits^n$, let
  $Q$ be a finite state machine of size $2^w$, and let $k\ge0$. Then
  \[
    \Prb[(h_1,\ldots,h_k)\text{ is not }((2^k-1)\varepsilon,Q)\text{-good}]
    < \frac{2^{6w}k}{\varepsilon^2 2^{2n}}.
  \]
\end{lemma}
\begin{proof}
  \uses{lem:good-induction-base, lem:bad-event-one, lem:bad-event-two, lem:induction-matrix-comparison, lem:induction-square-comparison}
  The case $k=0$ is Lemma~\ref{lem:good-induction-base}. For $k>0$, expose
  $h_1,\ldots,h_{k-1}$ and $h_k$. Event one is that the first $k-1$ hashes are
  good; its failure probability is bounded by Lemma~\ref{lem:bad-event-one}.
  Event two is the simultaneous independence of the last hash; its failure
  probability is bounded by Lemma~\ref{lem:bad-event-two}. By the union bound,
  both events fail with total probability at most
  $2^{6w}k/(\varepsilon^2 2^{2n})$. If both events hold, the triangle
  inequality plus Lemmas~\ref{lem:induction-matrix-comparison} and
  \ref{lem:induction-square-comparison} give
  \[
    \norm{Q(G_k(*))-Q((\Un_n)^{2^k})}_1
    \le \varepsilon+(2^k-2)\varepsilon=(2^k-1)\varepsilon.
  \]
  Therefore the complement of the good event has the stated probability.
\end{proof}

\begin{lemma}[Acceptance probability is controlled by matrix norm]
  \label{lem:acceptance-norm-bound}
  \uses{def:q-distribution, def:one-norms, def:good-sequence}
  If $(h_1,\ldots,h_k)$ is $(\eta,Q)$-good, then the difference between the
  acceptance probability of $Q$ on a uniform string in $(\bits^n)^{2^k}$ and
  on $G_k(x,h_1,\ldots,h_k)$ with uniform $x$ is at most $\eta$.
\end{lemma}
\begin{proof}
  \uses{def:q-distribution, def:one-norms, def:good-sequence}
  Let $e_{\start}$ be the row vector concentrated on the start state and let
  $1_{\accept}$ be the column vector that is $1$ on accepting states and $0$
  elsewhere. The two acceptance probabilities differ by
  \[
    e_{\start}\big(Q(G_k(*))-Q((\Un_n)^{2^k})\big)1_{\accept}.
  \]
  The vector $e_{\start}$ has one-norm $1$, and every entry of
  $1_{\accept}$ has absolute value at most $1$. The induced matrix norm
  therefore bounds the absolute value of the expression by $\eta$.
\end{proof}

\begin{lemma}[Recursive generator fools block-space machines; Lemma 3]
  \label{lem:generator-block-prg}
  \uses{def:block-prg, constr:recursive-generator, lem:good-sequence-probability, lem:acceptance-norm-bound}
  There is a constant $c>0$ such that for all integers $n$ and $k\le cn$, the
  recursive generator $G_k:\bits^n\times\calH^k\to(\bits^n)^{2^k}$ is a
  pseudorandom generator for $\Space(cn)$ and block size $n$ with parameter
  $2^{-cn}$.
\end{lemma}
\begin{proof}
  \uses{def:block-prg, lem:good-sequence-probability, lem:acceptance-norm-bound}
  Fix a $2^{cn}$-state machine $Q$. The acceptance gap is at most the
  probability that the sampled hash sequence is not good, plus the worst
  acceptance gap conditioned on goodness. Lemma~\ref{lem:acceptance-norm-bound}
  bounds the second term by the goodness parameter. Lemma~\ref{lem:good-sequence-probability}
  bounds the first term by $2^{2k}2^{6cn}k/(\varepsilon^2 2^{2n})$ after
  choosing the good parameter so that $(2^k-1)\varepsilon$ is the final norm
  error. Taking, for example, $c=0.05$ and the paper's
  $\varepsilon=2^{-cn-1}$ makes the sum at most $2^{-cn}$ for every
  $k\le cn$. This is precisely Definition~\ref{def:block-prg}.
\end{proof}

\chapter{Space-bounded derandomization and traversal sequences}

\begin{definition}[Pseudorandom generator for space-bounded algorithms; Definition 5]
  \label{def:space-prg}
  A map $G:\bits^m\to\bits^R$ is a pseudorandom generator for $\Space(S)$ with
  parameter $\varepsilon$ if for every randomized $\Space(S)$ algorithm $A$ and
  every fixed input to $A$,
  \[
    \abs{\Prb_{y\leftarrow\bits^R}[A(y)\text{ accepts}]
      -\Prb_{x\leftarrow\bits^m}[A(G(x))\text{ accepts}]}
    <\varepsilon .
  \]
\end{definition}

\begin{proposition}[Block generators are ordinary space generators; Proposition 1]
  \label{prop:block-to-space}
  \uses{def:block-fsm, def:block-prg, def:space-prg}
  If $G:\bits^m\to(\bits^n)^k$ is a pseudorandom generator for $\Space(S)$ and
  block size $n$ with parameter $\varepsilon$, then the concatenated output of
  $G$ is a pseudorandom generator for ordinary $\Space(S)$ algorithms with the
  same parameter $\varepsilon$.
\end{proposition}
\begin{proof}
  \uses{def:block-fsm, def:block-prg, def:space-prg}
  Fix a space-$S$ algorithm $A$ and a fixed input. Build the finite state
  machine $Q$ whose states are the configurations of $A$ between blocks of
  $n$ random bits. Its edge labels record exactly which $n$-bit random block
  moves one configuration to another. The number of configurations is at most
  $2^S$. Acceptance of $Q$ on a block sequence is the same event as acceptance
  of $A$ on the concatenated random string. Applying Definition~\ref{def:block-prg}
  to this $Q$ gives Definition~\ref{def:space-prg}.
\end{proof}

\begin{theorem}[Space-bounded generator; Theorem 1]
  \label{thm:space-prg-main}
  \uses{lem:generator-block-prg, prop:block-to-space, lem:convolution-universal}
  For all $R$ and $S$, there is an explicit pseudorandom generator
  \[
    G:\bits^{O(S\log R)}\to\bits^R
  \]
  for $\Space(S)$ with parameter $2^{-S}$. The generator can be computed in
  time polynomial in $R$ and $S$ and in space $O(S\log R)$.
\end{theorem}
\begin{proof}
  \uses{lem:generator-block-prg, prop:block-to-space, lem:convolution-universal}
  Choose the block length $\Theta(S)$ so that Lemma~\ref{lem:generator-block-prg}
  applies with error $2^{-\Theta(S)}$. Choose $k=\lceil\log R\rceil$ so the
  recursive generator outputs at least $R$ bits after concatenating its
  $2^k$ blocks, and discard any excess bits. Each hash function is represented
  with $O(S)$ bits by the convolution family of
  Lemma~\ref{lem:convolution-universal}; the seed consists of one
  $O(S)$-bit initial string plus $k$ such hash descriptions, for total length
  $O(S\log R)$. Proposition~\ref{prop:block-to-space} converts the block
  guarantee into the ordinary space-bounded guarantee. The recursive evaluation
  performs one explicit hash computation for each generated block and stores
  the current recursion data and hash descriptions, which fits in
  $O(S\log R)$ space.
\end{proof}

\begin{corollary}[Randomness reduction for randomized space algorithms; Corollary 1]
  \label{cor:randomness-reduction}
  \uses{thm:space-prg-main}
  Any randomized algorithm running in space $S$ and using $R$ random bits can
  be simulated by an algorithm using only $O(S\log R)$ random bits and
  $O(S\log R)$ space.
\end{corollary}
\begin{proof}
  \uses{thm:space-prg-main}
  Replace the original random tape by the output of the generator from
  Theorem~\ref{thm:space-prg-main}. The simulator samples only the generator
  seed, computes requested pseudorandom bits as needed, and runs the original
  algorithm. The distinguishing guarantee preserves the acceptance probability
  up to the theorem's parameter. Storing the seed and generator workspace costs
  $O(S\log R)$ space.
\end{proof}

\begin{corollary}[Black-box randomized Savitch simulation]
  \label{cor:black-box-savitch}
  \uses{thm:space-prg-main}
  Every $\RSPACE(S)$ computation using $R$ random bits can be simulated
  deterministically by enumerating all seeds of the Nisan generator, using
  $O(S\log R)$ space and the same simulation rule for every such randomized
  algorithm.
\end{corollary}
\begin{proof}
  \uses{thm:space-prg-main}
  Enumerate every seed of the generator in Theorem~\ref{thm:space-prg-main}
  and run the randomized machine using the corresponding generated random
  string. The enumeration uses only the seed counter and the generator
  workspace. Since the generator fools every space-$S$ algorithm, averaging
  over the seed enumeration approximates the original randomized acceptance
  probability. The procedure is black-box because it depends only on the
  machine's space bound, not on its internal transition structure.
\end{proof}

\begin{definition}[Universal traversal sequence]
  \label{def:universal-traversal-sequence}
  For $d$-regular graphs on $n$ vertices, a sequence over $\{1,\ldots,d\}$ is
  universal if, for every connected $d$-regular graph with a fixed ordering of
  the incident edges at each vertex and every starting vertex, following the
  edge indices in the sequence visits every vertex.
\end{definition}

\begin{lemma}[Logspace generators yield traversal sequences]
  \label{lem:prg-to-uts}
  \uses{def:space-prg, def:universal-traversal-sequence}
  Let $G$ be a pseudorandom generator that fools logarithmic-space machines
  with sufficiently small constant error and outputs strings encoding walks of
  length $L$ in $d$-regular $n$-vertex graphs. Then concatenating all possible
  outputs of $G$ gives a universal traversal sequence whose length is
  $L\cdot 2^{\text{seed length}}$.
\end{lemma}
\begin{proof}
  \uses{def:space-prg, def:universal-traversal-sequence}
  If the concatenation were not universal, then there would be a graph,
  port-numbering, and start vertex for which no generator output visits all
  vertices. A logarithmic-space tester can store the current vertex, the target
  vertex, and a step counter, and accept exactly those random walk strings that
  visit the target. Under truly random choices, a sufficiently long random
  walk has positive probability of reaching every fixed target in a connected
  regular graph. Under the generator, that probability would be zero for one
  target, contradicting the generator's small distinguishing error. Hence every
  target is reached by at least one generated walk, and the concatenation of
  all generated walks traverses the graph from every start.
\end{proof}

\begin{theorem}[Explicit universal traversal sequences; Theorem 2]
  \label{thm:universal-traversal}
  \uses{thm:space-prg-main, lem:prg-to-uts}
  For all $n$ and $2\le d<n$, there are explicit universal traversal sequences
  of length $n^{O(\log n)}$ for $d$-regular $n$-vertex graphs. Moreover, the
  sequences can be produced by a deterministic Turing machine using space
  logarithmic in the sequence length.
\end{theorem}
\begin{proof}
  \uses{thm:space-prg-main, lem:prg-to-uts}
  A walk step in a $d$-regular graph is encoded by $O(\log d)\le O(\log n)$
  random bits, and the logspace traversal tester of Lemma~\ref{lem:prg-to-uts}
  uses $O(\log n)$ space. Apply Theorem~\ref{thm:space-prg-main} with
  $S=O(\log n)$ and with polynomially many random bits for the walk length
  required by the standard random-walk traversal argument. The seed length is
  $O(\log^2 n)$, so enumerating all seeds multiplies the walk length by
  $2^{O(\log^2 n)}=n^{O(\log n)}$. The deterministic generator enumerator keeps
  a seed counter, a recursion stack, and the current output position, all of
  size logarithmic in the final sequence length.
\end{proof}

\chapter{Pseudo-independent block generation and amplification}

\begin{definition}[Pseudo-independent block generator; Definition 6]
  \label{def:pibg}
  A map $G:\bits^m\to(\bits^n)^k$ is a pseudo-independent block generator with
  parameter $\varepsilon$ if for every sequence of sets
  $A_1,\ldots,A_k\subseteq\bits^n$,
  \[
    \abs{\Prb[y_1\in A_1,\ldots,y_k\in A_k]-p_1\cdots p_k}\le\varepsilon,
  \]
  where $(y_1,\ldots,y_k)=G(x)$ for uniform $x$ and
  $p_i=|A_i|/2^n$.
\end{definition}

\begin{proposition}[Small-space block PRGs are pseudo-independent; Proposition 2]
  \label{prop:block-prg-to-pibg}
  \uses{def:block-prg, def:pibg}
  If $G:\bits^m\to(\bits^n)^k$ is a pseudorandom generator for
  $\Space(\log(k+2))$ and block size $n$ with parameter $\varepsilon$, then
  $G$ is a pseudo-independent block generator with parameter $\varepsilon$.
\end{proposition}
\begin{proof}
  \uses{def:block-prg, def:pibg}
  Given $A_1,\ldots,A_k$, build a finite state machine with states
  $0,1,\ldots,k$ and an extra fail state. The start state is $0$, the only
  accepting state is $k$, the edge from $i-1$ to $i$ is labelled by $A_i$, and
  the edge from $i-1$ to $\fail$ is labelled by $\bits^n\setminus A_i$.
  The machine has $k+2$ states, hence uses $\log(k+2)$ space. Under uniform
  independent blocks its acceptance probability is $p_1\cdots p_k$; under the
  generator it is exactly the left probability in Definition~\ref{def:pibg}.
  The block-PRG guarantee gives the claimed error.
\end{proof}

\begin{theorem}[Explicit pseudo-independent block generator; Theorem 3]
  \label{thm:pibg-main}
  \uses{lem:generator-block-prg, prop:block-prg-to-pibg}
  There is a constant $c>0$ such that for all integers $R$ and $k\le 2^R$,
  there is an explicit pseudo-independent block generator with parameter
  $2^{-R}$ that converts $cR\log k$ random bits into $k$ strings of length $R$.
\end{theorem}
\begin{proof}
  \uses{lem:generator-block-prg, prop:block-prg-to-pibg}
  Apply Lemma~\ref{lem:generator-block-prg} with block length $\Theta(R)$,
  space at least $\log(k+2)$, and enough recursive levels to output at least
  $k$ blocks. Since $k\le2^R$, the required number of levels is $O(\log k)$
  and lies in the permitted range after adjusting the absolute constant. If
  the block length is larger than $R$, discard excess bits from each block.
  Proposition~\ref{prop:block-prg-to-pibg} converts the block-space guarantee
  into pseudo-independence, and the seed consists of $O(\log k)$ hash
  descriptions of size $O(R)$.
\end{proof}

\begin{corollary}[Deterministic amplification]
  \label{cor:deterministic-amplification}
  \uses{thm:pibg-main}
  Let a one-sided-error randomized algorithm use $R$ random bits and succeed on
  yes-instances with probability at least $1/2$. For every $k\le R$, one can
  run $k$ amplified trials using $O(R\log k)$ random bits and reduce the
  probability that all trials fail to $(1+o(1))2^{-k}$. More generally, if the
  success probability is at least $1/\poly(k)$, polynomially many trials reduce
  failure to $2^{-k}$ while still using $O(R\log k)$ random bits.
\end{corollary}
\begin{proof}
  \uses{thm:pibg-main}
  Let $A_i$ be the set of random strings on which the $i$-th trial fails. For
  independent trials, the probability that all $k$ trials fail is the product
  of their individual failure probabilities, at most $2^{-k}$ in the
  probability-$1/2$ case. Theorem~\ref{thm:pibg-main} makes the joint failure
  probability differ from this product by at most $2^{-R}$, which is lower
  order for $k\le R$. If each trial succeeds with probability
  $1/\poly(k)$, take $\poly(k)$ trials so that the independent failure product
  is at most $2^{-k}$, and apply the same pseudo-independence estimate.
\end{proof}

\begin{proposition}[Two-sided amplification by majority]
  \label{prop:two-sided-amplification}
  \uses{thm:pibg-main}
  The pseudo-independent block generator also amplifies BPP-type algorithms:
  if each trial is correct with probability at least $2/3$, then using
  $O(R\log k)$ random bits for $k$ generated trials and taking the majority
  reduces the error exponentially in $k$, up to the generator error.
\end{proposition}
\begin{proof}
  \uses{thm:pibg-main}
  For each subset $T$ of more than $k/2$ trials, let $A_i$ be the bad-randomness
  set for trial $i$ if $i\in T$ and the full space otherwise. Pseudo-independence
  bounds the probability that all trials in $T$ are bad by the corresponding
  independent product plus $2^{-R}$. Summing over all majority-size subsets
  gives the usual Chernoff-style exponential upper bound, with the accumulated
  additive error controlled by choosing the theorem's parameter smaller than
  the target final error. Thus majority voting behaves as it would under
  independent randomness, up to the pseudorandom error.
\end{proof}

\chapter{Staged algorithms and random primes}

\begin{proposition}[Staged algorithm principle]
  \label{prop:staged-algorithm-principle}
  \uses{def:block-prg}
  Suppose a randomized computation is divided into stages, each stage consumes
  at most $n$ random bits, and the information retained between stages uses at
  most $w$ bits. Then replacing the stage randomness by the blocks of a
  pseudorandom generator for $\Space(w)$ and block size $n$ changes the final
  output distribution by at most the generator parameter, for every event that
  can be tested from the final state.
\end{proposition}
\begin{proof}
  \uses{def:block-fsm, def:block-prg}
  Model the retained inter-stage information as the state of a finite state
  machine, with one transition per stage labelled by the random block values
  that lead to each next retained state. The event of interest is represented
  by the accepting states after the final stage. Definition~\ref{def:block-prg}
  is exactly the statement that the acceptance probability changes by at most
  the generator parameter.
\end{proof}

\begin{algorithm}[Standard random-prime generator]
  \label{alg:uniform-prime-generation}
  Given $N=2^n$, repeat the following loop until success: choose a random
  integer $x$ in $\{1,\ldots,N\}$, test whether $x$ is prime, and stop if the
  test says prime. Output the final $x$.
\end{algorithm}

\begin{lemma}[Basic Rabin--Miller test stage]
  \label{lem:rabin-miller-basic-test}
  \uses{alg:uniform-prime-generation}
  There is a randomized primality-test stage using $O(n)$ random bits and
  $O(n)$ retained space such that every prime is accepted and every composite
  $n$-bit integer is rejected with probability at least $1/2$ by one stage.
  Repeating $O(n)$ independent stages makes the probability of accepting a
  composite exponentially small in $n$.
\end{lemma}
\begin{proof}
  \uses{alg:uniform-prime-generation}
  The Rabin--Miller test chooses a random witness modulo the candidate integer
  and performs modular exponentiations using $O(n)$-bit arithmetic. For a
  composite input, at least half of the possible witnesses reveal compositeness;
  for a prime input, no witness falsely rejects. The current candidate, the
  witness, and the modular-exponentiation workspace use $O(n)$ bits. Repeating
  the basic test $O(n)$ times multiplies the probability of missing
  compositeness by at most $2^{-O(n)}$.
\end{proof}

\begin{lemma}[The prime-search loop has linear expected length]
  \label{lem:prime-density-loop}
  \uses{alg:uniform-prime-generation}
  For $N=2^n$, the expected number of candidate choices made by
  Algorithm~\ref{alg:uniform-prime-generation} before a prime is found is
  $O(n)$.
\end{lemma}
\begin{proof}
  \uses{alg:uniform-prime-generation}
  The prime number theorem gives density $\Theta(1/\log N)=\Theta(1/n)$ for
  primes in $\{1,\ldots,N\}$. Each independently chosen candidate is prime
  with probability $\Theta(1/n)$, so the waiting time for the first prime is a
  geometric random variable with expectation $O(n)$.
\end{proof}

\begin{theorem}[Generating random $n$-bit primes with $O(n\log n)$ random bits]
  \label{thm:random-prime-low-randomness}
  \uses{prop:staged-algorithm-principle, thm:space-prg-main, alg:uniform-prime-generation, lem:rabin-miller-basic-test, lem:prime-density-loop}
  A nearly uniform random prime in the range $\{1,\ldots,2^n\}$ can be
  generated using only $O(n\log n)$ truly random bits.
\end{theorem}
\begin{proof}
  \uses{prop:staged-algorithm-principle, thm:space-prg-main, alg:uniform-prime-generation, lem:rabin-miller-basic-test, lem:prime-density-loop}
  Break Algorithm~\ref{alg:uniform-prime-generation} into stages: one stage
  chooses the candidate $x$, and each basic Rabin--Miller witness check is one
  additional stage. By Lemma~\ref{lem:rabin-miller-basic-test}, every stage
  consumes $O(n)$ random bits and retains only $O(n)$ bits between stages,
  mainly the candidate $x$ and loop counters. By Lemma~\ref{lem:prime-density-loop},
  $O(n)$ candidate stages suffice in expectation, and the primality checks can
  be repeated enough times that false acceptance of a composite has
  exponentially small probability. Apply Theorem~\ref{thm:space-prg-main}, or
  equivalently the staged form in Proposition~\ref{prop:staged-algorithm-principle},
  with space and block size $O(n)$ and polynomially many stages. The resulting
  seed length is $O(n\log n)$, and the pseudorandom execution changes the
  output distribution by only the chosen exponentially small parameter.
\end{proof}
